\subsection{Complex Numbers and Absolute Value}

\begin{definition} \label{def:complex-number-field-operations} \hlt{Complex Numbers}\\
A \hlt{complex number} is an ordered pair $(a,b)$ of real numbers. For $x=(a,b)$ and $y=(c,d)$, define equality componentwise and set $x+y\defeq(a+c,b+d)$ and $xy\defeq(ac-bd,ad+bc)$
\end{definition}

\begin{theorem} \label{thm:complex-numbers-form-field} \hlt{The Complex Numbers Form a Field}\\
With the operations in Definition \ref{def:complex-number-field-operations}, the set of complex numbers is a field.\\
Its additive and multiplicative identities are $(0,0)$ and $(1,0)$.
\begin{proof}
The field axioms are checked directly from the two coordinate formulas and the field structure of $\R$.\\
If $x=(a,b)\neq(0,0)$, then $a^2+b^2>0$ by Proposition \ref{prop:ordered-field-inequality-rules-analysis}, and $x^{-1}=\left(\frac{a}{a^2+b^2},\frac{-b}{a^2+b^2}\right)$ satisfies $xx^{-1}=(1,0)$.\\
\hlt{Structural Commentary}: Inverses use conjugate-over-square.\\
\hlt{Key Insight}: Positivity of $a^2+b^2$ supplies the denominator for every nonzero complex number.
\end{proof}
\end{theorem}

\begin{proposition} \label{prop:real-field-embedded-in-complex-field} \hlt{Real Numbers Inside Complex Numbers}\\
For all real $a,b$, $(a,0)+(b,0)=(a+b,0)$, and $(a,0)(b,0)=(ab,0)$.\\
Thus complex numbers of form $(a,0)$ have same arithmetic as real numbers, and identify $(a,0)$ with $a$.\\
The \hlt{complex unit} is then $i\defeq(0,1)$.
\end{proposition}

\begin{proposition} \label{prop:complex-unit-and-coordinate-form} \hlt{Coordinate Form and $i^2$}\\
We have $i^2=-1$. Moreover, if $a,b\in\R$, then $(a,b)=a+bi$.
\begin{proof}
First, $i^2=(0,1)(0,1)=(-1,0)=-1$. Also, $a+bi=(a,0)+(b,0)(0,1)=(a,0)+(0,b)=(a,b)$.\\
\hlt{Structural Commentary}: The ordered-pair construction recovers the usual notation after $i=(0,1)$ is named.\\
\hlt{Key Insight}: The symbol $i$ is not assumed; it is produced by the pair model.
\end{proof}
\end{proposition}

\begin{definition} \label{def:complex-conjugate-real-imaginary-parts} \hlt{Conjugate, Real Part, and Imaginary Part}\\
If $z=a+bi$ with $a,b\in\R$, then the \hlt{complex conjugate} of $z$ is $\overline{z}\defeq a-bi$.\\
The numbers $a$ and $b$ are the \hlt{real part} and \hlt{imaginary part}, written $a=\operatorname{Re}(z)$ and $b=\operatorname{Im}(z)$.
\end{definition}

\begin{proposition} \label{prop:complex-conjugation-identities} \hlt{Basic Conjugation Identities}\\
If $z,w\in\C$, then
\begin{enumerate}[label=(\roman*)]
\setlength{\itemsep}{0pt}
\item $\overline{z+w}=\overline{z}+\overline{w}$;
\item $\overline{zw}=\overline{z}\,\overline{w}$;
\item $z+\overline{z}=2\operatorname{Re}(z)$ and $z-\overline{z}=2i\operatorname{Im}(z)$;
\item $z\overline{z}$ is real and positive unless $z=0$.
\end{enumerate}
\begin{proof}
Write $z=a+bi$ and $w=c+di$.
\begin{enumerate}[label=(\roman*)]
\setlength{\itemsep}{0pt}
\item $\overline{z+w}=(a+c)-(b+d)i=\overline{z}+\overline{w}$.
\item $\overline{zw}=(ac-bd)-(ad+bc)i=(a-bi)(c-di)=\overline{z}\,\overline{w}$.
\item $z+\overline{z}=2a=2\operatorname{Re}(z)$ and $z-\overline{z}=2bi=2i\operatorname{Im}(z)$.
\item $z\overline{z}=a^2+b^2$, positive unless $a=b=0$, i.e. unless $z=0$.
\end{enumerate}
\hlt{Structural Commentary}: Conjugation fixes $\operatorname{Re}$ and negates $\operatorname{Im}$.\\
\hlt{Key Insight}: $z\overline{z}=a^2+b^2\geq0$.
\end{proof}
\end{proposition}

\begin{definition} \label{def:complex-absolute-value-analysis} \hlt{Absolute Value of a Complex Number}\\
For $z\in\C$, define $\abs{z}\defeq(z\overline{z})^{1/2}$, the nonnegative square root supplied by Theorem \ref{thm:nth-roots-positive-reals} and Proposition \ref{prop:complex-conjugation-identities}.
\end{definition}

\begin{theorem} \label{thm:complex-absolute-value-identities} \hlt{Absolute Value Identities}\\
For $z,w\in\C$,
\begin{enumerate}[label=(\roman*)]
\setlength{\itemsep}{0pt}
\item $\abs{z}>0$ unless $z=0$, and $\abs{0}=0$;
\item $\abs{\overline{z}}=\abs{z}$;
\item $\abs{zw}=\abs{z}\abs{w}$;
\item $\abs{\operatorname{Re}(z)}\leq\abs{z}$;
\item $\abs{z+w}\leq\abs{z}+\abs{w}$.
\end{enumerate}
\begin{proof}
\begin{enumerate}[label=(\roman*)]
\setlength{\itemsep}{0pt}
\item If $z=a+bi$, then $\abs{z}=(a^2+b^2)^{1/2}$; this is positive unless $a=b=0$, and $\abs{0}=0$.
\item $\abs{\overline{z}}^2=\overline{z}z=z\overline{z}=\abs{z}^2$, so nonnegative square roots give $\abs{\overline{z}}=\abs{z}$.
\item $\abs{zw}^2=zw\overline{zw}=z\overline{z}w\overline{w}=\abs{z}^2\abs{w}^2$; take nonnegative square roots.
\item If $z=a+bi$, then $\abs{\operatorname{Re}(z)}=\abs{a}\leq(a^2+b^2)^{1/2}=\abs{z}$.
\item $z\overline{w}$ and $\overline{z}w$ are conjugates, $\abs{z+w}^2 =\abs{z}^2+2\operatorname{Re}(z\overline{w})+\abs{w}^2 \leq\abs{z}^2+2\abs{z\overline{w}}+\abs{w}^2 =(\abs{z}+\abs{w})^2$, using (iii) and (iv). Take nonnegative square roots.
\end{enumerate}
\hlt{Structural Commentary}: The complex absolute value behaves like length because $z\overline{z}$ is a squared length.\\
\hlt{Key Insight}: Triangle inequality is a real-part estimate plus multiplicativity.
\end{proof}
\end{theorem}

\begin{theorem} \label{thm:finite-schwarz-inequality-complex-scalars} \hlt{Finite Schwarz Inequality for Complex Scalars}\\
If $a_1,\ldots,a_n$ and $b_1,\ldots,b_n$ are complex numbers, then $\abs{\sum\limits_{j=1}^{n}a_j\overline{b_j}}^2\leq \left(\sum\limits_{j=1}^{n}\abs{a_j}^2\right) \left(\sum\limits_{j=1}^{n}\abs{b_j}^2\right)$.
\begin{proof}
This is Theorem \ref{thm:cauchy-schwarz-inequality} in $\C^n$ with the standard inner product.\\
Equivalently, set $A=\sum\abs{a_j}^2$, $B=\sum\abs{b_j}^2$, and $C=\sum a_j\overline{b_j}$; when $B>0$, $\sum\abs{Ba_j-Cb_j}^2=B(AB-\abs{C}^2)\geq0$.\\
\hlt{Structural Commentary}: This scalar inequality is the coordinate model for geometric Cauchy--Schwarz inequality.\\
\hlt{Key Insight}: A nonnegative square sum controls the mixed term.
\end{proof}
\end{theorem}
